\chapter{التحليل المركب الموجّه لنظرية الأعداد}

\section{نطاق الفصل ومخرجاته}

\noindent\textbf{حالة الفصل:} \textenglish{VERIFIED}. النتائج الأساسية
موسومة بمصادرها أو براهينها، ونواة بيرون مؤجلة صراحة، ونجحت فحوص الجودة
والبناء. لا تشمل هذه البوابة المراجعة الثانية المستقلة.

يفترض الفصل معرفة التفاضل والتكامل والمتسلسلات والمستوى العقدي. هدفه
تثبيت الحزمة الدنيا التي ستستعمل لاحقًا، لا إحلاله محل مقرر كامل في
التحليل المركب. بنهايته ينبغي أن يطبق القارئ صيغة كوشي والبواقي ومبدأ
الحجة، ويميز بين مجال تقارب تمثيلٍ ما ومجال استمرار الدالة، ويعدّد
الشروط التي يجب فحصها قبل نقل مسار. المرجع القياسي للنتائج المقتبسة هو
\cite{Ahlfors1979}، ولتطبيقها على زيتا انظر \cite{Titchmarsh1986}.

\section{مقدمة منهجية}

لا تحتاج نظرية الأعداد التحليلية إلى كل موضوعات التحليل المركب بالدرجة نفسها. فالهدف في هذا الفصل ليس إعادة بناء نظرية الدوال العقدية كاملة، بل استخراج الأدوات التي تتكرر في دراسة دوال زيتا و\(L\)، والصيغ الصريحة، وتقدير المجاميع الحسابية.

سنركز على:

\begin{enumerate}
  \item الدوال الهولومورفية والميرومورفية.
  \item التكاملات المسارية وصيغة كوشي.
  \item الأصفار والأقطاب والبواقي.
  \item مبدأ الحجة ومبرهنة روشيه.
  \item متسلسلات لوران والاستمرار التحليلي.
  \item تقديرات كوشي ومبدأ القيمة العظمى.
  \item تحويل المسارات واستخراج الحدود الرئيسية من الأقطاب.
\end{enumerate}

\section{المستوى العقدي والدوال العقدية}

نكتب
\[
s=\sigma+it,
\qquad \sigma,t\in\mathbb{R}.
\]
في نظرية الأعداد التحليلية، يستخدم الرمز \(s\) كثيرًا لأن سلاسل ديريشليه تكتب عادة على الصورة
\[
F(s)=\sum_{n=1}^{\infty}\frac{a_n}{n^s}.
\]

إذا كان \(n>0\)، فإن
\[
n^{-s}=e^{-s\log n}=n^{-\sigma}e^{-it\log n}.
\]
ومن هنا يظهر فصل مهم بين:

\begin{itemize}
  \item الحجم، الذي تحكمه \(n^{-\sigma}\).
  \item التذبذب، الذي يحكمه \(e^{-it\log n}\).
\end{itemize}

ولهذا فإن الخطوط الرأسية \(\sigma=\text{ثابت}\) تلعب دورًا مركزيًا في دراسة سلاسل ديريشليه.

\section{الاشتقاق العقدي ومعادلات كوشي--ريمان}
\symbolindex{الدالة الهولومورفية}
\symbolindex{معادلات كوشي--ريمان}

لتكن \(f:\Omega\to\mathbb{C}\)، حيث \(\Omega\subseteq\mathbb{C}\) مفتوحة. نقول إن \(f\) قابلة للاشتقاق عقديًا عند \(z_0\) إذا كانت النهاية
\[
f'(z_0)=\lim_{z\to z_0}\frac{f(z)-f(z_0)}{z-z_0}
\]
موجودة ولا تعتمد على اتجاه الاقتراب.

إذا كتبنا
\[
f(x+iy)=u(x,y)+iv(x,y),
\]
فإن القابلية للاشتقاق العقدي تفرض، تحت فرضيات انتظام مناسبة، معادلتي كوشي--ريمان:
\[
\frac{\partial u}{\partial x}
=
\frac{\partial v}{\partial y},
\qquad
\frac{\partial u}{\partial y}
=
-\frac{\partial v}{\partial x}.
\]

\begin{definition}
تسمى \(f\) هولومورفية على \(\Omega\) إذا كانت قابلة للاشتقاق عقديًا عند كل نقطة من \(\Omega\).
\end{definition}

القابلية للاشتقاق العقدي أقوى بكثير من القابلية للاشتقاق الحقيقي. فالدالة الهولومورفية قابلة للاشتقاق عددًا غير منتهٍ من المرات، بل تمثل محليًا بمتسلسلة قوى.

\section{متسلسلات القوى}

إذا كانت
\[
\sum_{n=0}^{\infty}a_n(z-z_0)^n
\]
ذات نصف قطر تقارب \(R>0\)، فإنها تتقارب تقاربًا مطلقًا ومحليًا منتظمًا داخل القرص \(|z-z_0|<R\)، وتمثل دالة هولومورفية هناك.

كما يمكن الاشتقاق حدًا حدًا:
\[
\left(\sum_{n=0}^{\infty}a_n(z-z_0)^n\right)'
=
\sum_{n=1}^{\infty}na_n(z-z_0)^{n-1}.
\]

\begin{remark}
التقارب المحلي المنتظم هو الجسر الصحيح بين التحليل الحدّي والهولومورفية. أما التقارب النقطي وحده فلا يكفي عادة للحفاظ على الخواص التحليلية.
\end{remark}

\section{المسارات والتكامل العقدي}

ليكن \(\gamma:[a,b]\to\mathbb{C}\) مسارًا أملسًا على أجزاء. نعرف
\[
\int_\gamma f(z)\,dz
=
\int_a^b f(\gamma(t))\gamma'(t)\,dt.
\]

إذا عكسنا اتجاه المسار، تغيرت إشارة التكامل. وإذا وصلنا مسارين، كان التكامل على المسار المركب مجموع التكاملين.

\begin{definition}
إذا كان \(\gamma\) مغلقًا، أي \(\gamma(a)=\gamma(b)\)، نكتب أحيانًا
\[
\oint_\gamma f(z)\,dz.
\]
\end{definition}

في التطبيقات العددية، تكون المسارات غالبًا مستطيلات أو خطوطًا رأسية وأفقية، لأن تغيير خط التكامل يسمح بالتقاط بواقي الأقطاب.

\section{مبرهنة كوشي}
\personindex{كوشي، أوغستان لوي}

\begin{theorem}[كوشي]
\theoremindex{مبرهنة كوشي}
\label{thm:cauchy-integral-theorem}
\resultid{ANT-THM-03-01}
\citedresult{\cite{Ahlfors1979}}
إذا كانت \(f\) هولومورفية على مجال بسيط الاتصال \(\Omega\)، وكان \(\gamma\) مسارًا مغلقًا داخل \(\Omega\)، فإن
\[
\oint_\gamma f(z)\,dz=0.
\]
\end{theorem}

تكمن قوة هذه المبرهنة في أن التكامل لا يعتمد على الشكل التفصيلي للمسار ما دام التشوه لا يعبر مفردات.

\section{صيغة كوشي التكاملية}

\begin{theorem}[صيغة كوشي]
\theoremindex{صيغة كوشي التكاملية}
\label{thm:cauchy-integral-formula}
\resultid{ANT-THM-03-04}
\citedresult{\cite{Ahlfors1979}}
إذا كانت \(f\) هولومورفية على جوار قرص مغلق مركزه \(z_0\) ونصف قطره \(R\)، وكان \(|z-z_0|<R\)، فإن
\[
f(z)
=
\frac{1}{2\pi i}
\oint_{|w-z_0|=R}
\frac{f(w)}{w-z}\,dw.
\]
\end{theorem}

وعلى وجه الخصوص،
\[
f^{(n)}(z)
=
\frac{n!}{2\pi i}
\oint_{|w-z_0|=R}
\frac{f(w)}{(w-z)^{n+1}}\,dw.
\]

\section{تقديرات كوشي}

\begin{theorem}[تقديرات كوشي]
\theoremindex{تقديرات كوشي}
\label{thm:cauchy-estimates}
\resultid{ANT-THM-03-05}
\provedhere
إذا كانت
\[
|f(w)|\leq M
\]
على الدائرة \(|w-z_0|=R\)، فإن
\[
|f^{(n)}(z_0)|
\leq
\frac{n!M}{R^n}.
\]
\end{theorem}

\begin{proof}
نطبق صيغة كوشي للمشتقة:
\[
f^{(n)}(z_0)
=
\frac{n!}{2\pi i}
\oint_{|w-z_0|=R}
\frac{f(w)}{(w-z_0)^{n+1}}\,dw.
\]
طول الدائرة \(2\pi R\)، ومن ثم
\[
|f^{(n)}(z_0)|
\leq
\frac{n!}{2\pi}
\cdot 2\pi R
\cdot \frac{M}{R^{n+1}}
=
\frac{n!M}{R^n}.
\]
\end{proof}

هذه التقديرات مهمة عند التحكم في مشتقات دوال زيتا و\(L\)، أو عند تحويل حدود النمو على دائرة إلى معلومات محلية.

\section{مبدأ القيمة العظمى}

\begin{theorem}[مبدأ القيمة العظمى]
\theoremindex{مبدأ القيمة العظمى}
\label{thm:maximum-modulus}
\resultid{ANT-THM-03-06}
\citedresult{\cite{Ahlfors1979}}
إذا كانت \(f\) هولومورفية وغير ثابتة على مجال متصل، فلا يمكن أن تحقق \(|f|\) قيمة عظمى محلية داخل المجال.
\end{theorem}

ومن نتائجه:

\begin{itemize}
  \item دالة هولومورفية محدودة على المستوى كله ثابتة، وهي مبرهنة ليوفيل.
  \item يمكن نقل تقديرات الحدود إلى داخل المجالات.
  \item تظهر مبادئ من نمط فراجمن--ليندلوف عند التعامل مع مجالات غير محدودة.
\end{itemize}

\section{الأصفار ورتبتها}

إذا كانت \(f\) هولومورفية قرب \(z_0\)، وكان
\[
f(z_0)=0,
\]
فإما أن \(f\) منعدمة تطابقًا على المركبة المتصلة، أو يوجد عدد صحيح \(m\geq1\) ودالة هولومورفية \(g\) لا تنعدم عند \(z_0\) بحيث
\[
f(z)=(z-z_0)^m g(z).
\]

يسمى \(m\) رتبة الصفر.

\begin{corollary}
\label{cor:isolated-zeros}
\resultid{ANT-COR-03-01}
\citedresult{\cite{Ahlfors1979}}
أصفار الدالة الهولومورفية غير المنعدمة تطابقًا معزولة.
\end{corollary}

هذه الحقيقة أساسية لأن أصفار دوال زيتا و\(L\) يمكن عدها داخل مناطق محددة.

\section{المفردات المعزولة}

إذا كانت \(f\) هولومورفية في جوار مثقوب لـ\(z_0\)، فإنها تملك توسع لوران:
\[
f(z)=\sum_{n=-\infty}^{\infty}a_n(z-z_0)^n.
\]

تصنف المفردة إلى:

\begin{enumerate}
  \item مفردة قابلة للإزالة إذا كان \(a_n=0\) لكل \(n<0\).
  \item قطب من الرتبة \(m\) إذا كان \(a_{-m}\neq0\) ولا توجد حدود سالبة أدنى منه.
  \item مفردة جوهرية إذا ظهر عدد لا نهائي من الحدود السالبة.
\end{enumerate}

\begin{definition}
الدالة الميرومورفية هي دالة هولومورفية عدا أقطاب معزولة.
\end{definition}

دالة زيتا لريمان بعد الاستمرار التحليلي ميرومورفية على المستوى كله،
ولها قطب بسيط عند \(s=1\) \cite{Titchmarsh1986}.

\section{الباقي}

\begin{definition}
باقي \(f\) عند \(z_0\) هو معامل \((z-z_0)^{-1}\) في توسع لوران:
\[
\operatorname{Res}_{z=z_0}f(z)=a_{-1}.
\]
\end{definition}

إذا كان \(z_0\) قطبًا بسيطًا، فإن
\[
\operatorname{Res}_{z=z_0}f(z)
=
\lim_{z\to z_0}(z-z_0)f(z).
\]

وإذا كان قطبًا من الرتبة \(m\)، فإن
\[
\operatorname{Res}_{z=z_0}f(z)
=
\frac{1}{(m-1)!}
\lim_{z\to z_0}
\frac{d^{m-1}}{dz^{m-1}}
\left((z-z_0)^m f(z)\right).
\]

\section{مبرهنة البواقي}

\begin{theorem}[مبرهنة البواقي]
\label{thm:residue-theorem}
\resultid{ANT-THM-03-02}
\citedresult{\cite{Ahlfors1979}}
إذا كانت \(f\) ميرومورفية داخل مسار مغلق \(\gamma\) وعلى جواره، ولا توجد أقطاب على المسار، فإن
\[
\oint_\gamma f(z)\,dz
=
2\pi i
\sum_{\rho}
\operatorname{Res}_{z=\rho}f(z),
\]
حيث يمتد المجموع على الأقطاب داخل المسار مع مراعاة عدد اللف.
\end{theorem}

هذه المبرهنة هي الآلة التي تحول الأقطاب إلى حدود رئيسية. ففي نظرية الأعداد، يؤدي نقل خط تكامل عبر قطب إلى ظهور مساهمة من باقي ذلك القطب.

\section{مثال: استخراج حد رئيسي من قطب بسيط}

افترض أن \(F(s)\) ميرومورفية قرب \(s=1\)، ولها قطب بسيط:
\[
F(s)=\frac{A}{s-1}+H(s),
\]
حيث \(H\) هولومورفية.

إذا ظهر التكامل
\[
\frac{1}{2\pi i}\int_{c-iT}^{c+iT}F(s)\frac{x^s}{s}\,ds
\]
مع \(c>1\)، ثم نقلنا الخط إلى يسار \(1\)، فإن القطب عند \(s=1\) يساهم بالمقدار
\[
\operatorname{Res}_{s=1}
\left(F(s)\frac{x^s}{s}\right)
=
Ax.
\]

هذه الصورة المجردة تفسر ظهور الحدود الرئيسية الخطية في كثير من مسائل العد.

\section{مبدأ الحجة}

\begin{theorem}[مبدأ الحجة]
\label{thm:argument-principle}
\resultid{ANT-THM-03-03}
\provedhere
إذا كانت \(f\) ميرومورفية داخل \(\gamma\) وعلى جواره، ولا تملك أصفارًا أو أقطابًا على \(\gamma\)، فإن
\[
\frac{1}{2\pi i}
\oint_\gamma
\frac{f'(z)}{f(z)}\,dz
=
N-P,
\]
حيث \(N\) عدد الأصفار و\(P\) عدد الأقطاب داخل المسار، مع حساب الرتب.
\end{theorem}

\begin{proof}
إذا كان \(\rho\) صفرًا من الرتبة \(m\)، نكتب
\[
f(z)=(z-\rho)^m g(z),
\]
حيث \(g(\rho)\neq0\). ومن ثم
\[
\frac{f'}{f}(z)
=
\frac{m}{z-\rho}+\frac{g'}{g}(z),
\]
فيكون باقي \(f'/f\) عند \(\rho\) مساويًا \(m\). وبالمثل، القطب من الرتبة \(m\) يعطي باقيًا \(-m\). ثم نطبق مبرهنة البواقي.
\end{proof}

هذه الصيغة هي الأساس في عد أصفار دوال زيتا و\(L\).

\section{مبرهنة روشيه}

\begin{theorem}[روشيه]
\label{thm:rouche}
\resultid{ANT-THM-03-07}
\citedresult{\cite{Ahlfors1979}}
لتكن \(f\) و\(g\) هولومورفيتين داخل مسار بسيط مغلق وعلى جواره. إذا
\[
|g(z)|<|f(z)|
\]
على المسار، فإن \(f\) و\(f+g\) لهما العدد نفسه من الأصفار داخل المسار، مع حساب الرتب.
\end{theorem}

المعنى العملي: إذا كانت \(g\) اضطرابًا أصغر من \(f\) على الحدود، فلا يغير عدد الأصفار داخل المنطقة.

\section{الاستمرار التحليلي}

\begin{definition}
إذا كانت \(f\) هولومورفية على مجال \(\Omega_1\)، و\(F\) هولومورفية على مجال أكبر \(\Omega_2\supseteq\Omega_1\)، وكان
\[
F|_{\Omega_1}=f,
\]
فإن \(F\) استمرار تحليلي لـ\(f\).
\end{definition}

\begin{theorem}[مبدأ الهوية]
\label{thm:identity-theorem}
\resultid{ANT-THM-03-08}
\citedresult{\cite{Ahlfors1979}}
إذا اتفقت دالتان هولومورفيتان على مجموعة لها نقطة تراكم داخل مجال متصل، فإنهما متطابقتان على المجال كله.
\end{theorem}

لذلك يكون الاستمرار التحليلي، إذا وجد، وحيدًا.

\section{مثال أولي للاستمرار التحليلي}

السلسلة
\[
\sum_{n=0}^{\infty}z^n
\]
تتقارب فقط عندما \(|z|<1\)، لكنها تساوي هناك
\[
\frac{1}{1-z}.
\]
والدالة \(1/(1-z)\) معرفة ميرومورفيًا على المستوى عدا \(z=1\). إذن التعبير الكسري استمرار تحليلي لمجموع السلسلة خارج قرص التقارب.

هذه الفكرة تتكرر مع دالة زيتا: السلسلة الأصلية تتقارب في نصف مستوى، لكن الدالة نفسها تمتد إلى مجال أكبر.

\section{تحويل المسارات}

إذا كانت \(f\) هولومورفية بين مسارين رأسيين، يمكن نقل التكامل من أحدهما إلى الآخر دون تغيير، بعد إضافة القطع الأفقية وإثبات أن مساهمتها صغيرة أو تزول.

أما إذا عبرنا أقطابًا، فإن الفرق بين التكاملين يساوي مجموع البواقي المضروبة في \(2\pi i\).

\begin{remark}
الخطوة الأصعب غالبًا ليست كتابة مبرهنة البواقي، بل تبرير:
\begin{enumerate}
  \item تقارب التكامل.
  \item زوال القطع الأفقية.
  \item عدم عبور أصفار أو أقطاب غير محسوبة.
  \item الانتظام في المعلمات.
\end{enumerate}
\end{remark}

\section{تكاملات ميلين}

تحويل ميلين لدالة \(f\) هو
\[
\mathcal{M}f(s)
=
\int_0^\infty f(x)x^{s-1}\,dx.
\]

يحوّل هذا التحويل التدرج الضربي إلى تحليل عقدي في المتغير \(s\). وهو نظير لتحويل فورييه بعد التغيير \(x=e^u\).

في نظرية الأعداد، تتفاعل تحويلات ميلين طبيعيًا مع سلاسل ديريشليه، وتظهر في:

\begin{itemize}
  \item الصيغ العكسية.
  \item النوى الملساء.
  \item المعادلات الوظيفية.
  \item تحليل الدوال الخاصة.
\end{itemize}

\section{تمهيد لصيغة بيرون}

إذا كانت
\[
F(s)=\sum_{n=1}^{\infty}\frac{a_n}{n^s}
\]
متقاربة مطلقًا عندما \(\Re(s)>c_0\)، فإن صيغة بيرون تربط مجموع المعاملات
\[
\sum_{n\leq x}a_n
\]
بتكامل رأسي من الشكل
\[
\frac{1}{2\pi i}
\int_{c-iT}^{c+iT}
F(s)\frac{x^s}{s}\,ds.
\]

جوهر الصيغة هو هوية الانقلاب الحدّية الآتية.

\begin{theorem}[نواة بيرون]
\label{thm:perron-kernel}
\resultid{ANT-THM-03-09}
\deferredresult{الفصل المخصص لصيغة بيرون}
لكل \(c>0\) و\(y>0\)، تُفهم النهاية المتناظرة على أنها
\[
\lim_{T\to\infty}
\frac{1}{2\pi i}
\int_{c-iT}^{c+iT}
\frac{y^s}{s}\,ds
=
\begin{cases}
1,& y>1,\\
\frac12,& y=1,\\
0,& 0<y<1.
\end{cases}
\]
\end{theorem}

هذا التكامل ليس مطلق التقارب؛ ولذلك لا يجوز استعمال الصيغة غير المقطوعة
من دون تحديد معنى النهاية وضبط إمكان تبديلها مع سلسلة ديريشليه. لم
تُثبت النتيجة بعد داخل الموسوعة، وستُثبت مع نسختي بيرون المقطوعة
والملساء وفروضهما في فصل مستقل؛ انظر مرجعًا \cite{Titchmarsh1986}.

\section{أخطاء شائعة}

\begin{enumerate}
  \item نقل المسار دون فحص النمو على القطع الأفقية.
  \item استعمال مبرهنة البواقي مع وجود قطب على المسار.
  \item الخلط بين نصف قطر تقارب السلسلة ومجال الاستمرار التحليلي للدالة.
  \item نسيان رتب الأصفار والأقطاب في مبدأ الحجة.
  \item افتراض أن التقارب النقطي يحفظ الهولومورفية.
  \item استعمال اللوغاريتم العقدي دون تحديد الفرع.
\end{enumerate}

\section{تمارين}

\begin{enumerate}
  \item أثبت أن الدالة الهولومورفية التي تنعدم على متتالية لها نقطة تراكم داخل المجال منعدمة تطابقًا.
  \item احسب باقي
  \[
  \frac{e^z}{(z-1)^2}
  \]
  عند \(z=1\).
  \item استعمل مبرهنة البواقي لحساب
  \[
  \oint_{|z|=2}\frac{dz}{z(z-1)}.
  \]
  \item أثبت أن
  \[
  \frac{f'}{f}
  \]
  له باقي يساوي رتبة الصفر عند كل صفر لـ\(f\).
  \item طبق روشيه لإثبات أن كثير الحدود
  \[
  z^5+10z+1
  \]
  له عدد محدد من الأصفار داخل دائرة مناسبة.
  \item أثبت تقديرات كوشي للمشتقات.
  \item فسّر كيف ينتج الحد \(Ax\) من قطب بسيط ذي باقي \(A\) في تكامل يحتوي على \(x^s/s\).
\end{enumerate}

\section{خلاصة الفصل}

التحليل المركب في نظرية الأعداد ليس مجرد خلفية تقنية؛ بل هو آلة لتحويل البنية التحليلية إلى معلومات حسابية. فالأقطاب تنتج حدودًا رئيسية، والأصفار تنتج تذبذبًا، ومبدأ الحجة يعد الأصفار، وتحويل المسارات ينقل التكامل إلى مناطق أكثر فائدة، والاستمرار التحليلي يكشف معلومات تتجاوز مجال تقارب السلسلة الأصلية.
